\section{MoE 稳定性、fine-tuning 与 upcycling}

前面的 routing 与 systems 章节说明 MoE 能扩展容量，但训练结束并不代表模型可稳定使用。本节解释 router z-loss 如何抑制极端 logits、fine-tuning 为什么可能让少数 experts 过拟合，以及 dense-to-MoE upcycling 如何复用已有权重，降低从头训练大规模稀疏模型的成本。


\begin{figure}[H]
\centering
\includegraphics[width=0.88\textwidth]{slide-047.jpg}
\caption{Slide 48: Issues with MoEs - stability}
\figsource{CS336 Spring 2026 Lecture 4 官方幻灯片。}
\end{figure}

\noindent\textbf{展开说明：}[Zoph 2022]

\begin{knowledgebox}{读图：Slide 48 应该怎么看}
这页是机制或证据页。先看它比较的是 compute、参数量、routing、负载均衡还是通信；再看结论支持的是质量提升、训练速度、推理成本还是系统复杂度。MoE 和 attention alternatives 的图常把模型结构和系统代价混在一起，读图时要分清“算法机制”和“硬件执行成本”。
\end{knowledgebox}


\begin{figure}[H]
\centering
\includegraphics[width=0.88\textwidth]{slide-048.jpg}
\caption{Slide 49: Z-loss stability for the router}
\figsource{CS336 Spring 2026 Lecture 4 官方幻灯片。}
\end{figure}

\noindent\textbf{展开说明：}What happens when we remove the z-loss?

\begin{knowledgebox}{读图：Slide 49 应该怎么看}
这页是机制或证据页。先看它比较的是 compute、参数量、routing、负载均衡还是通信；再看结论支持的是质量提升、训练速度、推理成本还是系统复杂度。MoE 和 attention alternatives 的图常把模型结构和系统代价混在一起，读图时要分清“算法机制”和“硬件执行成本”。
\end{knowledgebox}


\begin{figure}[H]
\centering
\includegraphics[width=0.88\textwidth]{slide-049.jpg}
\caption{Slide 50: Issues with MoEs – fine-tuning}
\figsource{CS336 Spring 2026 Lecture 4 官方幻灯片。}
\end{figure}

\noindent\textbf{展开说明：}Sparse MoEs can overfit on smaller fine-tuning data

\begin{knowledgebox}{读图：Slide 50 应该怎么看}
这页是机制或证据页。先看它比较的是 compute、参数量、routing、负载均衡还是通信；再看结论支持的是质量提升、训练速度、推理成本还是系统复杂度。MoE 和 attention alternatives 的图常把模型结构和系统代价混在一起，读图时要分清“算法机制”和“硬件执行成本”。
\end{knowledgebox}


\begin{figure}[H]
\centering
\includegraphics[width=0.88\textwidth]{slide-050.jpg}
\caption{Slide 51: Other training methods - upcycling}
\figsource{CS336 Spring 2026 Lecture 4 官方幻灯片。}
\end{figure}

\noindent\textbf{展开说明：}Can we use a pre-trained LM to initialize a MoE?

\begin{knowledgebox}{读图：Slide 51 应该怎么看}
这页是机制或证据页。先看它比较的是 compute、参数量、routing、负载均衡还是通信；再看结论支持的是质量提升、训练速度、推理成本还是系统复杂度。MoE 和 attention alternatives 的图常把模型结构和系统代价混在一起，读图时要分清“算法机制”和“硬件执行成本”。
\end{knowledgebox}



\begin{figure}[H]
\centering
\includegraphics[width=0.88\textwidth]{slide-051.jpg}
\caption{Slide 52：MiniCPM 的 dense-to-MoE upcycling 案例。}
\figsource{CS336 Spring 2026 Lecture 4 官方幻灯片。}
\end{figure}

\noindent\textbf{展开说明：}Uses the MiniCPM model (topk=2, 8 experts, ~ 4B active params). Upcycling 的关键不是简单复制 FFN，而是让新 experts 初始行为接近 dense checkpoint、router 不立即坍缩，并在后续训练中逐渐分化。它节省从随机初始化开始的成本，但也可能复制原模型偏差，且总存储与通信会立刻上升。


\begin{figure}[H]
\centering
\includegraphics[width=0.88\textwidth]{slide-052.jpg}
\caption{Slide 53: Upcycling example – Qwen MoE}
\figsource{CS336 Spring 2026 Lecture 4 官方幻灯片。}
\end{figure}

\noindent\textbf{展开说明：}Qwen MoE – Initialized from the Qwen 1.8B model top-k=4, 60 experts w/ 4 shared.

\subsection{本章小结}
本节的共同线索是 selective computation：通过结构选择让 token 不必总是访问所有历史或所有参数。但选择性越强，越需要额外机制保证信息流、负载均衡和训练稳定。

